\section{Contexto y antecedentes}

Ethereum se ha consolidado como la plataforma de referencia para el despliegue de contratos
inteligentes y constituye la base sobre la que se sostiene una porción significativa del
ecosistema de finanzas descentralizadas~\cite{buterin2014ethereum}. Este crecimiento no se
produjo de forma aislada, sino que estuvo acompañado por la proliferación de redes de
\textit{Layer 2} y cadenas alternativas que replican o extienden su modelo de ejecución, cada
una optimizada según distintos criterios de costo, rendimiento o descentralización. La
fragmentación resultante de este ecosistema plantea la necesidad de que los activos digitales y
los datos puedan trasladarse entre las distintas cadenas que lo componen, requisito que se
conoce como interoperabilidad. Los puentes entre cadenas, o \textit{cross-chain bridges},
constituyen la infraestructura que hace posible dicho traslado, de modo que la seguridad de
estos mecanismos determina, en gran medida, la seguridad de los fondos que a través de ellos se
transfieren.

Paralelamente al desarrollo de estas redes, las pruebas de conocimiento cero
(\textit{Zero-Knowledge Proofs}) evolucionaron desde ser un objeto de estudio predominantemente
teórico hasta convertirse en una herramienta de ingeniería aplicable a problemas concretos. Esta
tecnología permite comprimir la verificación de un cómputo arbitrariamente extenso en una prueba
sucinta, cuya validación puede realizarse en una fracción del tiempo y del costo que exigiría
reejecutar dicho cómputo. Aplicada al contexto de Ethereum, esta propiedad habilita la
verificación liviana, desde otra cadena, tanto del consenso alcanzado por los validadores —a
través de la firma emitida por el comité de sincronización que certifica la finalidad de un
bloque— como de la inclusión efectiva de una transacción particular dentro de ese bloque, lo
cual se logra mediante una prueba sobre su árbol de Merkle-Patricia. La combinación de ambas
verificaciones, comprimidas en una única prueba de conocimiento cero, constituye el fundamento
técnico sobre el que resulta posible construir un puente entre cadenas cuya seguridad no dependa
de la honestidad de un intermediario, sino de una demostración matemáticamente verificable. En
el presente trabajo, este fundamento se materializa mediante el diseño e implementación de un
puente del tipo \textit{lock-and-mint}, en el cual el bloqueo de un activo en la cadena de
origen habilita su acuñación equivalente en la cadena de destino únicamente cuando dicha prueba
ha sido validada \textit{on-chain}.

\section{Planteamiento del problema}

La verificación directa del consenso de Ethereum y de la inclusión de una transacción, sin
comprimir dicho cómputo mediante una prueba de conocimiento cero, resulta impracticable cuando
se pretende realizar fuera de la propia cadena. Reproducir la agregación de las firmas del
comité de sincronización y recorrer un árbol de Merkle-Patricia dentro de un contrato
inteligente desplegado en otra cadena implica un volumen considerable de operaciones
criptográficas, fundamentalmente sobre curvas elípticas, cuyo costo en gas resulta, en la
práctica, prohibitivo. Esta limitación conduce a que los sistemas de interoperabilidad
existentes adopten un compromiso distinto, en el cual la verificación criptográfica del estado
de la cadena de origen es sustituida por la delegación de dicha verificación en un comité de
validadores externos sobre el cual se deposita confianza. Este diseño traslada el problema de
seguridad del propio protocolo hacia el comportamiento honesto de un conjunto reducido de
actores, y esta delegación de confianza ha constituido, de manera sistemática, el punto de
falla explotado en los incidentes de seguridad más costosos sufridos por puentes entre cadenas
hasta la fecha~\cite{zhang2023bridgesok}. Persiste, en consecuencia, una tensión no resuelta
entre dos objetivos igualmente necesarios: por un lado, alcanzar garantías criptográficas
equivalentes a las del consenso de la cadena de origen y, por otro, mantener un costo de
verificación \textit{on-chain}, tanto en gas como en tiempo de generación de la prueba, que
resulte viable para un sistema de uso real.

\section{Justificación}

El ecosistema de herramientas orientadas al desarrollo de circuitos de conocimiento cero
capaces de abordar este volumen de cómputo criptográfico se encuentra, a su vez, en un proceso
activo de maduración, y ofrece actualmente dos paradigmas de diseño claramente diferenciados.
El primero está representado por los lenguajes de dominio específico, orientados
exclusivamente a la escritura de circuitos, entre los cuales se encuentra Noir~\cite{noirlang},
que expone al desarrollador primitivas y un modelo de programación concebidos desde su diseño
para la generación eficiente de pruebas. El segundo paradigma corresponde a las máquinas
virtuales de conocimiento cero de propósito general, cuyo exponente más reciente es
SP1~\cite{succinctsp1}, las cuales permiten probar la ejecución de programas escritos en un
lenguaje de propósito general, en este caso Rust, compilados a un conjunto de instrucciones que
la máquina virtual ejecuta y prueba sin que el desarrollador deba razonar directamente en
términos de restricciones aritméticas. Si bien ambos paradigmas persiguen resolver el mismo
tipo de problema, lo hacen mediante compromisos de diseño distintos entre la expresividad del
lenguaje, la facilidad de desarrollo y depuración, y el rendimiento final del circuito
resultante. Dado que ambas herramientas son de adopción reciente y se encuentran en desarrollo
activo, no existe todavía evidencia empírica consolidada que permita fundamentar la elección de
uno u otro paradigma para la implementación de un caso de uso concreto y computacionalmente
exigente, como lo es la verificación de consenso e inclusión de transacciones en Ethereum. El
presente trabajo se propone generar dicha evidencia mediante la implementación de un mismo
problema en ambas tecnologías, bajo condiciones comparables, y la medición sistemática de su
comportamiento. Ambas implementaciones se integran, además, dentro del mismo puente
\textit{lock-and-mint}, de manera que la comparación entre Noir y SP1 no se limita a un
ejercicio aislado de generación de pruebas, sino que se valida en el contexto de una
transferencia real de activos digitales entre cadenas.

\section{Objetivos}

\subsection{Objetivo general}

Diseñar, implementar y comparar empíricamente la verificación de consenso de Ethereum y de
inclusión de transacciones mediante Noir y SP1, integrando ambas implementaciones en un puente
\textit{lock-and-mint} de confianza reducida, con el propósito de determinar cuál de ambas
tecnologías ofrece un mejor rendimiento, un menor costo de verificación y una mayor viabilidad
práctica para la transferencia segura de activos digitales entre cadenas.

\subsection{Objetivos específicos}

\begin{itemize}
	\item Investigar los mecanismos de prueba y las arquitecturas de circuitos que subyacen a
	Noir, en tanto lenguaje de dominio específico, y a SP1, en tanto máquina virtual de
	conocimiento cero de propósito general basada en Rust.
	\item Implementar, en ambas tecnologías, la lógica de verificación de encabezados de bloque
	de Ethereum, de la firma del comité de sincronización y de las pruebas de inclusión de
	transacciones.
	\item Medir y comparar las métricas de rendimiento relevantes entre ambas implementaciones,
	entre las que se incluyen el tiempo de generación de la prueba, el uso de recursos
	computacionales y de memoria, el costo de verificación \textit{on-chain} en gas, y la
	complejidad relativa del proceso de desarrollo.
\end{itemize}

\section{Hipótesis}

El presente trabajo se orienta por la siguiente pregunta de investigación: ¿presenta el
paradigma de máquina virtual de conocimiento cero de propósito general, representado por SP1,
ventajas significativas en términos de flexibilidad de desarrollo y de tiempos de ejecución
frente a un lenguaje de conocimiento cero de dominio específico como Noir, al implementar la
verificación de consenso e inclusión de transacciones en Ethereum?

\section{Delimitación}

El presente trabajo se circunscribe a la arquitectura de consenso de Ethereum posterior a
\textit{The Merge}, particularmente al mecanismo del comité de sincronización empleado por los
clientes ligeros para certificar la finalidad de un bloque, así como a las pruebas de inclusión
sobre su árbol de Merkle-Patricia. Las implementaciones desarrolladas y las comparativas de
rendimiento realizadas se validan sobre redes de prueba, empleando Ethereum Sepolia como cadena
de origen y Optimism Sepolia como cadena de destino, bajo una infraestructura de cómputo
estandarizada que permite que las mediciones obtenidas para Noir y para SP1 resulten
comparables entre sí. Quedan fuera del alcance de este trabajo tanto la implementación sobre
redes principales como la verificación de mecanismos de consenso correspondientes a cadenas
distintas de Ethereum.

